\documentclass[12pt]{article}

\usepackage[utf8]{inputenc}
\usepackage[T1]{fontenc}
\usepackage[spanish,es-nodecimaldot]{babel}
\usepackage[margin=2.5cm]{geometry}
\usepackage{booktabs}
\usepackage{amsmath}
\usepackage[hidelinks]{hyperref}
\usepackage{parskip}
\usepackage{titlesec}
\usepackage{xcolor}
\usepackage{enumitem}
\usepackage{fancyhdr}
\usepackage{array}
\setlength{\headheight}{14.5pt}

\definecolor{impgreen}{RGB}{0,128,64}
\definecolor{darkgray}{RGB}{60,60,60}

\pagestyle{fancy}
\fancyhf{}
\fancyhead[L]{\small\textcolor{darkgray}{Reporte de Avances — Comité Tutorial, Junio 2026}}
\fancyhead[R]{\small\textcolor{darkgray}{F.~J. Cisneros Ruiz}}
\fancyfoot[C]{\small\thepage}
\renewcommand{\headrulewidth}{0.4pt}

\titleformat{\section}{\normalsize\bfseries\color{impgreen}}{}{0em}{}[\titlerule]
\titlespacing{\section}{0pt}{10pt}{4pt}

\begin{document}

% ── Portada ──────────────────────────────────────────────────────────────
\begin{titlepage}
  \centering
  \vspace*{1.5cm}
  {\large\bfseries Instituto Mexicano del Petróleo\\[0.3em]
  \normalsize Coordinación Académica de Posgrado — Doctorado en Ingeniería\par}

  \vspace{1.8cm}
  \rule{\textwidth}{1.2pt}\\[0.5em]
  {\Large\bfseries Reporte de Avances\\[0.2em]
  \normalsize Comité Tutorial — Último semestre\par}
  \rule{\textwidth}{1.2pt}

  \vspace{1.5cm}
  \begin{tabular}{ll}
    \textbf{Estudiante:}           & Francisco Javier Cisneros Ruiz \\[0.4em]
    \textbf{Matrícula:}            & 92735 \quad Generación 2022 \\[0.4em]
    \textbf{Director de tesis:}    & Dr.\ Iván Félix González \\[0.4em]
    \textbf{Semestre:}             & Enero -- Junio 2026 (último semestre regular) \\[0.4em]
    \textbf{Sede actual:}          & Veracruz, Veracruz \\[0.4em]
    \textbf{Fecha de evaluación:}  & 26 de junio de 2026, 12:40 h \\
  \end{tabular}

  \vspace{1.8cm}
  {\large\bfseries
  Detección de Daños en Plataformas Marinas tipo Jacket\\
  mediante Optimización Inversa con Algoritmos Genéticos\\
  e Índice de Calidad de Detección (ICD)}

  \vfill
  {\small Ciudad de México / Veracruz, 22 de junio de 2026}
\end{titlepage}

% ── Sección 1 ─────────────────────────────────────────────────────────────
\section{Contexto del Proyecto Doctoral}

Esta tesis propone una metodología de \textbf{Monitoreo de Salud Estructural (SHM)} basada en
\textbf{Algoritmos Genéticos (AG)} para identificar daños en plataformas \textit{Jacket} del Golfo de México.
La contribución central es el \textbf{Índice de Calidad de Detección (ICD)}, que fusiona ocho
indicadores de respuesta vibratoria y penaliza exponencialmente los falsos positivos.
La validación se realizó sobre \textbf{3\,240 corridas de simulación} bajo escenarios de corrosión y
abolladura (\textit{denting}) en un modelo de elementos finitos calibrado.

% ── Sección 2 ─────────────────────────────────────────────────────────────
\section{Avances del Semestre (Enero -- Junio 2026)}

\textbf{Cuatro avances principales} se consolidaron este semestre:

\medskip
\noindent\textbf{1.\enspace Procesador de Fatiga SACS — Producto de desarrollo (IMP)}

Como parte del proceso de titulación del Doctorado en Ingeniería del IMP, se completó la
\textbf{segunda de tres participaciones en proyecto institucional}: una herramienta de
escritorio desarrollada para los ingenieros estructurales del IMP
(coordinadora: Dra.\ Claudia Rendón Conde) que automatiza la consolidación de reportes
de análisis de fatiga generados por el software SACS. El problema que resuelve es concreto:
SACS genera un archivo de reporte independiente por cada período de operación analizado,
y su suma manual es propensa a errores. La herramienta carga $N$ archivos, extrae y normaliza
los datos, suma los daños por elemento y los clasifica en cuatro niveles de riesgo (CRÍTICO,
ALTO, MODERADO, BAJO), exportando un reporte Excel con formato profesional.

La interfaz gráfica (\textit{desktop GUI}, esquema oscuro, tres paneles) opera sin necesidad de
Python en el equipo del usuario y se distribuye como ejecutable Windows independiente (\texttt{.exe}).
Fue validada con datos reales de una plataforma marina: 506\,239 líneas de texto procesadas
sin errores, 2\,304 incidencias consolidadas en menos de 8 segundos.
La suite de pruebas automatizadas cuenta con \textbf{99 pruebas (100\,\% aprobadas)}.

\medskip
\noindent\textbf{2.\enspace Envío del artículo a \textit{Marine Structures} (Elsevier)}

El manuscrito doctoral fue enviado el \textbf{21 de junio de 2026} a la revista
\textit{Marine Structures} (Elsevier), cambiando el destino original (\textit{Applied Ocean Research})
debido a que esta última opera actualmente en modalidad exclusivamente \textit{Open Access},
con un costo de publicación (APC) inviable para el estudiante.
\textit{Marine Structures} mantiene una indexación y factor de impacto equivalentes sin costo para el autor.

\medskip
\noindent\textbf{3.\enspace Refinaciones del ICD y análisis de vectores de pesos ($\alpha_1$--$\alpha_8$)}

Se realizó un análisis estadístico profundo de los pesos que el AG asigna a cada uno de los
ocho indicadores de daño, permitiendo contextualizar mejor el ICD: identificar qué indicadores
dominan según el tipo de elemento (arriostres vs.\ piernas) y el mecanismo de daño,
y confirmar que el índice escala de forma diferenciada y controlada con la severidad.

\medskip
\noindent\textbf{4.\enspace Nueva métrica: Precisión Ponderada por Distancia Topológica}

Se propone $\text{Precisión}_w$, que asigna a cada falso positivo (FP) un peso
según su distancia $d$ en el grafo de conectividad estructural:
\[
  w(d) = \frac{d}{d+1}
\]
La Precisión clásica penaliza igual a un FP adyacente al nodo dañado que a uno en el extremo
opuesto de la estructura, subestimando la calidad diagnóstica real. $\text{Precisión}_w$
corrige esta limitación, con ganancias de hasta \textbf{$+$22 puntos porcentuales} en
elementos tipo \textit{brace}.

% ── Sección 3 ─────────────────────────────────────────────────────────────
\section{Participación en Proyecto IMP en Curso}

Como parte del proceso de titulación, el estudiante se encuentra actualmente realizando
su \textbf{tercera y última participación en proyecto institucional} del IMP, en colaboración
con el \textbf{Dr.\ Francisco Silva González} del Instituto Mexicano del Petróleo.
El proyecto tiene una duración estimada de \textbf{tres meses} y representa el cumplimiento
final de este requisito indispensable para la obtención del grado.

Paralelamente, se está llevando a cabo la \textbf{redacción del documento de tesis doctoral},
actividad cuya duración estimada es comparable a la del proyecto en curso.

% ── Sección 4 ─────────────────────────────────────────────────────────────
\section{Estado General y Solicitud de Prórroga}

\begin{table}[h!]
  \centering
  \small
  \caption{Estado de actividades del semestre.}
  \begin{tabular}{lc}
    \toprule
    \textbf{Actividad} & \textbf{Estado} \\
    \midrule
    Procesador de Fatiga SACS entregado al IMP         & \textcolor{impgreen}{\textbf{Completado}} \\
    Envío del artículo a \textit{Marine Structures}    & \textcolor{impgreen}{\textbf{Completado}} \\
    Refinaciones del ICD y análisis de vectores alpha  & \textcolor{impgreen}{\textbf{Completado}} \\
    Propuesta de Precisión Ponderada Topológica        & \textcolor{impgreen}{\textbf{Completado}} \\
    3ra participación IMP (Dr.\ Francisco Silva González) & \textcolor{orange}{\textbf{En curso}} \\
    Redacción del documento de tesis doctoral            & \textcolor{orange}{\textbf{En curso}} \\
    \bottomrule
  \end{tabular}
\end{table}

El presente semestre constituye el \textbf{último semestre regular} del programa doctoral.
Dado que el proceso de revisión por pares en revistas indexadas del área puede extenderse
entre seis meses y un año, se solicitará formalmente ante el Comité Tutorial una
\textbf{prórroga de un año} para dar tiempo a la aceptación del artículo, requisito
indispensable para la defensa de tesis.

Los dos pendientes en curso ---la tercera participación en proyecto IMP y la redacción
de la tesis doctoral--- concluirán \textbf{antes o en el mismo plazo que la aceptación
de la publicación en revista}, por lo que no representan un retraso adicional respecto
al horizonte de titulación.

Los detalles técnicos de cada avance se presentarán en la \textbf{última evaluación
de desempeño del 26 de junio de 2026 a las 12:40 h}.

\end{document}
